% Extracted from part_III_IV.tex for inclusion in main.tex
\setcounter{section}{3}

\section{Multi-Agent Reinforcement Learning Under Imperfect Monitoring}
\label{sec:marl}
% ============================================================

The analytical frameworks of Sections~2 and~3 share a common
simplifying premise: cooperation is either characterised under
common knowledge of primitives (demand intercept, cost vector,
rivals' strategies) or governed by a population-level adjustment
rule that operates outside the agents' own decision problem.
Neither abstraction corresponds to the empirical reality of
crude-oil markets, where producers observe a continuously quoted
benchmark price but receive official production figures with
substantial lags, frequent revisions, and non-trivial cross-agency
disagreement.\footnote{Monthly output estimates published by OPEC,
the International Energy Agency, and the U.S.\ Energy Information
Administration routinely diverge by several hundred thousand barrels
per day for the same reference month, with revisions arriving many
quarters after the fact. Brent, WTI, and Dubai benchmark prices,
by contrast, are quoted continuously and are available to all market
participants in real time.}

This section therefore asks a sharply delimited empirical question:
can model-free algorithmic agents, which receive the
realised market price as their sole feedback signal and never
observe rivals' individual production, spontaneously converge to
supra-competitive output levels? The question has acquired concrete
regulatory urgency over the past five years. \citet{calvano2020}
establish that two independent $\varepsilon$-greedy $Q$-learners
in a differentiated Bertrand game reliably sustain supra-Nash
pricing without any communication or programmed coordination.
Subsequent contributions have extended that observation to richer
architectures and sequential posted-price environments
\citep{klein2021, asker2022, banchio2023}. Regulators have
responded: California's AB~325, operative since January~2026,
criminalises algorithmic intermediaries that ``facilitate price
alignment'' \citep{dlapiper2025}; the European Commission has
confirmed multiple active investigations \citep{nortonrose2025,loyensloeff2026};
and the U.S.\ Department of Justice has publicly flagged the risk
of ``fully automated cartels'' \citep{kavanagh2025}. The open
empirical question motivating the present section is whether the
\citeauthor{calvano2020} mechanism, documented in differentiated
Bertrand competition, survives the translation to \emph{Cournot
quantity competition with asymmetric costs and imperfect
monitoring}, which is the institutional structure of OPEC+.

\paragraph{Why the theory of Sections~2-3 leaves the question
open.}
The Folk Theorem is a result about existence, not
selection. It guarantees that, for patient enough
producers, a collusive path can be supported as a
subgame-perfect equilibrium; it says nothing about which
of the continuum of sustainable equilibria will actually be
played. Classical analysis closes that indeterminacy by
assumption, that is, by endowing agents with common knowledge of
demand, costs, and rivals' strategies, and by letting them compute
the punishment thresholds that deter deviation. Learning agents
have none of these. They observe only the realised price, hold no
model of the demand system, cannot anticipate a rival's
best response, and revise behaviour by reinforcement rather than
by strategic reasoning. The economically interesting question is
therefore not whether cooperation is feasible, which the
theory already settles, but whether and how it is
selected by adaptive agents who cannot reason their way to
it. If cooperation emerges, the next question is whether it is
sustained by the Folk-Theorem mechanism (anticipated punishment)
or by something else entirely. Anticipating our central result:
cooperation does emerge, but the punishment mechanism the theory
predicts does not, and that gap is the substantive economic
contribution of this section.

Five primary hypotheses, formulated before any large-budget run was
initiated, organise the experimental programme. They serve as
public pre-commitments against post-hoc selection of decision
thresholds; all five are evaluated verbatim in the verdict matrix
of Section~\ref{subsec:matrix}.

\begin{enumerate}\itemsep0pt
  \item[\textbf{H1.}] \emph{Triopoly collusion.}~The
        greedy-rollout collusion index satisfies
        $\CI_{\text{trio}} > 0.30$
        (Section~\ref{subsec:headline}).
  \item[\textbf{H2.}] \emph{Learner-count ordering.}~
        $\CI_{\text{trio}} > \CI_{\text{duo}} >
        \CI_{\text{single}}$, with non-overlapping $95\%$
        confidence intervals
        (Section~\ref{subsec:learners}).
  \item[\textbf{H3.}] \emph{Folk-Theorem monotonicity.}~
        Spearman $\rho(\gamma, \CI) > 0$ on the patience
        sweep (Section~\ref{subsec:gamma}).
  \item[\textbf{H4.}] \emph{Green-Porter punishment cycles.}~
        Shock-driven retaliation phases are detected at
        $\geq 2\sigma_P$ in the headline run
        (Section~\ref{subsec:shocks}).
  \item[\textbf{H5.}] \emph{Stackelberg recovery.}~The
        leader-output learner reproduces the static leader
        quantity within $\pm 20\%$
        (Section~\ref{subsec:stack}).
\end{enumerate}

Departures from these hypotheses are reported as substantive
findings in their own right, not reframed to preserve a positive
narrative.

\begin{table}[H]
\centering
\caption*{\textbf{Results at a glance.} Pre-registered verdicts, each
developed in the section indicated and consolidated in the verdict
matrix (Table~\ref{tab:matrix}). H4's shock-response operationalisation
is met (\S\ref{subsec:shocks}); its three retaliation tests are
not (\S\ref{subsec:falsification}), which is the section's central
finding. \\[2pt] \badgelegend}
\small
\renewcommand{\arraystretch}{1.2}
\begin{tabular}{clcc}
\toprule
 & Pre-registered hypothesis & Verdict & Section \\
\midrule
\textbf{H1} & Triopoly collusion emerges & \vyes & \ref{subsec:headline} \\
\textbf{H2} & More learners $\Rightarrow$ more collusion & \vyes & \ref{subsec:learners} \\
\textbf{H3} & Collusion rises with patience $\gamma$ & \vyes & \ref{subsec:gamma} \\
\textbf{H4} & Cooperation held by Folk-Theorem punishment & \vno & \ref{subsec:falsification} \\
\textbf{H5} & Single learner recovers Stackelberg leader & \vpart & \ref{subsec:stack} \\
\bottomrule
\end{tabular}
\end{table}

% ============================================================
\subsection{Algorithm, information structure, and protocol}
\label{subsec:setup}
% ============================================================

\begin{notationbox}\small
$\CI$ -- collusion index (Eq.~\eqref{eq:ci}): $0$ at static Nash,
$1$ at the joint-profit cartel. \enspace
$\CI^{\text{tail}}$ vs.\ $\CI^{\text{greedy}}$ -- tail-of-training
($\varepsilon=0.05$) vs.\ frozen greedy-rollout ($\varepsilon=0$)
estimator. \enspace
$\sigma_P$ -- standard deviation of the rolling-mean price over the
converged window. \enspace
$\rho$ -- Spearman rank correlation (reciprocity tests). \enspace
$n_{\text{seeds}}$ -- independent random-seed replications. \enspace
$\gamma=\delta$ -- discount factor (algorithmic patience).
\end{notationbox}

\paragraph{Environment.}
The stage game is the linear-Cournot triopoly calibrated in
Section~\ref{sec:foundations}: inverse demand $P(Q) = a - bQ$ with $a = 140$, $b = 1$,
and marginal costs $c_{\text{US}} = 45$, $c_{\text{OPEC}} = 20$,
$c_{\text{RUS}} = 35$. Quantity action spaces are unconstrained
unless otherwise stated. The market clears every period; each agent
observes the common price and its own realised profit but not
rivals' individual quantities. This price-only information
structure is exactly the Green-Porter environment studied
analytically in Section~\ref{subsec:greenporter}. Here, however,
the agents are model-free learners who must discover the implicit
rules of the game through repeated interaction, rather than agents
endowed with a parametric model of the demand system.

\paragraph{Algorithm.}
Each agent runs an independent tabular $\varepsilon$-greedy
$Q$-learner \citep{watkins1992}\citep[p.~??]{sutton2018}:
\begin{equation}
  \label{eq:qlearn}
  Q_i(s_t, a_t) \;\leftarrow\; Q_i(s_t, a_t)
    \;+\; \alpha\bigl[\,r_{i,t}
      + \gamma\max_{a'} Q_i(s_{t+1}, a')
      - Q_i(s_t, a_t)\,\bigr],
\end{equation}
where $r_{i,t} = (P_t - c_i)\,q_{i,t}$ is the per-period profit,
$\alpha \in (0,1)$ is the learning rate, and $\gamma \in (0,1)$
is the discount factor, the algorithmic counterpart of the
patience parameter $\delta$ in the Folk-Theorem analysis of
Section~\ref{subsec:folk}. Agents are deliberately kept independent: no shared
value function, no communication channel, no centralised reward
signal.

\paragraph{State and action spaces.}
The private state $s_t^i = (\hat{P}_t, \hat{q}_t^{\,i})$ encodes
the last-period market price and the agent's own last-period
output, discretised into $12$ price bins and $10$ own-quantity
bins. Actions are drawn from a $15$-point grid spanning
$[0, q_i^{\max}]$, where $q_i^{\max}$ equals twice the
cooperative cartel quota, a range wide enough to accommodate both
joint-profit maximisation and outright market flooding. Exploration
follows a geometric annealing schedule:
$\varepsilon_t = \max\!\{\varepsilon_{\text{end}},\,
\varepsilon_{\text{start}}\cdot\rho^{\,t}\}$,
with $(\varepsilon_{\text{start}}, \varepsilon_{\text{end}},
\rho) = (0.30,\, 0.05,\, 0.99)$ and $\alpha = 0.15$.

\paragraph{Calibrated parameters versus algorithmic
hyperparameters.}
We draw a sharp distinction between two classes of inputs. The
economic parameters, that is, the demand intercept and
slope $(a,b)$, the marginal-cost vector, and the discount factor
$\delta=\gamma$, are calibrated to empirical sources
(\citealp{dallasfed2026,bp2022,imf2021}; see Section~\ref{sec:foundations}) and carry economic meaning that can be
disciplined by data. The algorithmic hyperparameters, that
is, the learning rate $\alpha$, the exploration schedule
$(\varepsilon_{\text{start}},\varepsilon_{\text{end}},\rho)$, and
the grid resolution, have no empirical counterpart: they describe
how an artificial agent learns, not how the oil market behaves.
We therefore make no claim to ``estimate'' them. Instead, we fix
each at its literature-standard default and sweep it across
the recommended range in the robustness battery of
Section~\ref{subsec:robustness}, reporting the headline finding as
robust precisely because it survives every value rather than
because any single value is correct. This is the appropriate
standard of justification for a parameter that is computational
rather than economic.

\paragraph{Outcome metric.}
The unifying metric across Sections~2-4 is the collusion index
\begin{equation}
  \label{eq:ci}
  \CI \;=\;
    \frac{\bar{P} - P^{\,\text{Nash}}}
         {P^{\,\text{cartel}} - P^{\,\text{Nash}}},
\end{equation}
where $\bar{P}$ is the mean price over the final $20\%$ of the
training horizon. By construction, $\CI = 0$ under the static
Cournot-Nash allocation, $\CI = 1$ under the joint-profit cartel,
and $\CI \in (0,1)$ on the segment of tacit cooperation. Values
below zero (encountered in the single-learner and duopoly regimes
documented below) identify predatory pricing below static Nash.

\paragraph{Aggregation.}
Every table and figure reports a multi-seed mean with either a
per-seed standard deviation (and the corresponding
normal-approximation $95\%$ confidence interval) for scalar
quantities, or a per-step $2.5\%$-$97.5\%$ percentile band for
trajectory quantities. The headline triopoly batch of
Section~\ref{subsec:headline} draws on $n_{\text{seeds}} = 50$;
the robustness battery of
Sections~\ref{subsec:gamma}-\ref{subsec:stack} uses
$n_{\text{seeds}} = 20$ at half the headline horizon
($750\,\text{episodes} \times 50\,\text{steps}$); the cycle and
reciprocity checks of Sections~\ref{subsec:cycles} and
\ref{subsec:reciprocity} use $n_{\text{seeds}} = 5$ at the full
headline budget.

% ============================================================
\subsection{Headline experiment: three independent Q-learners}
\label{subsec:headline}
% ============================================================

\emph{In brief:} across $50$ independent seeds of the calibrated
triopoly, three price-only learners collude without any coordination
device, the cooperative basin is reached on $\mathbf{94\%}$ of seeds
with a mean greedy collusion index of $\mathbf{0.73}$ ($95\%$ CI
$[0.62,\,0.84]$), clearing the pre-registered H1 threshold of $0.30$
(\textbf{H1 met}).

Under the Folk Theorem of Section~\ref{subsec:folk} and the Green-Porter
framework of Section~\ref{subsec:greenporter}, three sufficiently patient agents
observing the realised price can in principle sustain any
cooperative output above static Nash, provided the discount factor
exceeds the binding critical patience $\delta^\ast$ and imperfect
monitoring does not destroy cooperation in expectation. Tabular
$Q$-learning is the minimal model-free algorithm that, upon
convergence, can implement such a reciprocity rule. Whether it
actually discovers a cooperative basin from random
initialisation and price-only feedback is the first question the
experiment addresses.

The triopoly is trained from scratch for $1{,}500$ episodes of
$50$ steps each on $50$ independent random seeds. All
hyper-parameters are as in Section~\ref{subsec:setup}. Two
estimators of the converged collusion index are computed: a
\emph{tail-of-training} estimator (mean over the final $20\%$ of
training, which still carries $\varepsilon = 0.05$ exploration
noise) and a \emph{greedy rollout} estimator (Q-tables frozen,
$\varepsilon = 0$, re-evaluated on $20$ noise-free rollouts of
$50$ steps each). Pre-registered inference is based on the greedy
estimator, which is directly interpretable as the converged
policy's implied market outcome.

\begin{criterion*}
Tacit algorithmic cooperation is declared to emerge iff the
large-budget run satisfies all three of: \textup{(i)} the $95\%$
CI on the mean greedy $\CI$ is strictly above~$0$;
\textup{(ii)} its lower bound is $\geq 0.30$; and
\textup{(iii)} the per-seed standard deviation of $\bar{P}$ over
the converged window is $\leq 5\,\$\text{/bbl}$.
\end{criterion*}

Table~\ref{tab:headline} reports the central statistics of the
$50$-seed batch.

\begin{verdict*}
Criteria \textup{(i)} and \textup{(ii)} are satisfied: the $95\%$
CI on the mean greedy $\CI$ is $[0.62,\,0.84]$, strictly above
zero with a lower bound exceeding $0.30$. Criterion~\textup{(iii)}
is satisfied on the tail estimator
($\sigma_{\bar{P}} = 3.60 < 5\,\$\text{/bbl}$).
\textbf{H1~is~met.}
\end{verdict*}

Two features of the data merit closer examination before this
finding is interpreted. First, the random-policy
baseline (uniform draws on the $15$-level action grid) yields
$\CI = 0.04$ ($95\%$ CI $[-0.04,\,0.12]$), confirming
that the cooperative attractor is not a discretisation artefact of
the output grid. Second, $74\%$ of state-action cells are visited
fewer than five times. This high under-visitation rate is a
structural property of the discretisation rather than an artefact
of insufficient training; its implications for the mechanism
interpretation are taken up in
Section~\ref{subsec:reciprocity}.

% --- TABLE 1 ---

% ============================================================
\subsection{Per-seed distribution and multi-basin structure}
\label{subsec:basins}
% ============================================================

A cross-seed mean of $\CI^{\text{greedy}} = 0.73$
compresses fifty independent training histories into a single
number, and the compression is misleading. Figure~\ref{fig:basins}
makes the underlying distribution explicit: independent tabular
$Q$-learning does not converge to a unique equilibrium: it
converges to a distribution over qualitatively distinct
attractors. Table~\ref{tab:basins} classifies each seed by its
greedy collusion index into four basins.

This bimodality reframes the headline finding along two
dimensions. First, the tail-of-training estimator (cross-seed
standard deviation~$0.18$) systematically understates the
cross-seed dispersion of the greedy estimator ($0.39$), because
residual $\varepsilon$-noise smooths each agent's effective policy
across nearby $Q$-cells. Both estimators should therefore be
reported side by side; the gap between them is itself a diagnostic
of under-training.

Second, when the empirical distribution is bimodal, the
appropriate summary is not the cross-seed mean but the
proportion of seeds reaching a cooperative basin.
Applying a \citet{wilson1927} score interval, the standard
correction for normal-approximation breakdown at extreme
proportions, applied to $\Pr[\CI \geq 0.20]$ yields the comparison in
Table~\ref{tab:wilson}. The triopoly and duopoly Wilson intervals
do not overlap on either estimator (lower bound $0.84$ versus
upper bound $0.07$): with three learners, the probability of
reaching a cooperative basin is bounded below by $0.84$;
replacing a single learner with a rational myopic best-responder
collapses that probability below $0.07$.

The cooperative basin is therefore not a property of the action
grid, the static benchmark, or the hyper-parameter configuration.
It is caused by all players learning simultaneously. The
reformulated headline statement, which is robust to the choice of
estimator and invariant to the bimodality, reads:

\begin{quote}
  \itshape
  Under price-only information and independent tabular
  $Q$-learning, the probability that a triopoly of independent
  learners reaches a cooperative basin ($\CI \geq 0.20$) is
  $0.94$,\footnote{This cross-seed basin frequency is a distinct
  quantity from the within-run cooperative-regime time fraction
  (also $94\%$) reported for the Green-Porter AR(1) simulation in
  Section~\ref{subsec:greenporter}; the numerical coincidence is
  incidental.} with a Wilson $95\%$ interval of $[0.84,\,0.98]$. In
  a duopoly where one agent is replaced by a myopic
  best-responder, this probability is $0.00$
  ($[0.00,\,0.07]$).
\end{quote}

% --- TABLE 2 ---

% --- TABLE 3 ---

% ============================================================
\subsection{Duopoly robustness check}
\label{subsec:duopoly}
% ============================================================

\emph{In brief:} replacing one learner with a static Cournot
best-responder (OPEC and the US learning, Russia myopic) does not merely
weaken cooperation but flips the outcome to a predatory
equilibrium ($\CI = -0.21$, $95\%$ CI $[-0.24,\,-0.18]$), confirming
that the cooperative basin is a property of all-player learning rather
than of the environment.

To isolate the role of all-player learning, the headline
experiment is replicated with only two learning agents (OPEC
and the US) while Russia plays the static Cournot best-response
at every step. The pre-registered criterion requires the duopoly
$95\%$ CI on $\CI$ to sit strictly below the triopoly CI; overlap
would render the regime contrast undetectable.

The duopoly yields a mean collusion index of $-0.21 \pm 0.10$
($95\%$ CI $[-0.24,\,-0.18]$) and a mean converged price of
$55.81 \pm 1.98\,\$\text{/bbl}$ over $50$ seeds. Against the
triopoly headline ($\CI = 0.63$, $95\%$ CI $[0.57,\,0.67]$), the
two intervals are disjoint by more than $0.80$ units. The
pre-registered criterion is satisfied.

The substantive point is stronger than the interval comparison
suggests. The duopoly does not deliver a weaker form of
cooperation: it delivers a predatory equilibrium, with
prices roughly four dollars below static Nash. The presence of
even one myopic best-responder is sufficient to destroy the
cooperative basin entirely; a finding examined in further detail
in Section~\ref{subsec:learners}.

% ============================================================
\subsection{Falsification of explicit Folk-Theorem retaliation}
\label{subsec:falsification}
% ============================================================

\emph{In brief:} three independent tests of whether cooperation is
sustained by anticipated punishment, as the Folk Theorem requires forced-deviation decomposition, $\sigma$-calibrated cycle detection, and
Spearman reciprocity at fine and coarse grids, all reject explicit
retaliation, identifying the cooperative outcome as a path-dependent
lock-in rather than a punishment equilibrium.

The preceding sections establish that cooperation emerges and is
causally attributable to all-player learning. They do not,
however, identify the mechanism by which it is sustained. The
Friedman-Abreu taxonomy offers the canonical candidate: rational
cooperators punish deviators by reverting to Nash (grim trigger)
or to a temporary flooding phase \citep{abreu1986}. If the
cooperative basins documented above are algorithmic realisations
of such an equilibrium, three distinct empirical signatures should
be detectable: a non-trivial strategic share in the price
response to a forced deviation; \emph{$\sigma$-calibrated
punishment cycles} of detectable amplitude and frequency; and a
state-conditioned negative correlation between observed
prices and subsequent outputs. Each test is applied in turn.

% ------------------------------------------------------------
\subsubsection{Forced-deviation decomposition}
\label{subsec:forced}
% ------------------------------------------------------------

Pinning one learner's output at a Nash-like level for a fixed
window is the algorithmic analogue of the grim-trigger experiment
of Section~\ref{subsec:folk}. Two distinct price effects are predicted. A
mechanical drop $\Delta P_{\text{mech}} = b\,\Delta
q_{\text{dev}}$ follows directly from the deviator's additional
supply meeting the inverse-demand slope; this is not evidence of
punishment. A strategic drop $\Delta P_{\text{strategic}}$,
arising from non-deviators raising their own outputs in response,
is the distinctive signature of algorithmic retaliation.

OPEC's quantity is fixed at $q^{\text{forced}} = 40\,\text{mbd}$
for $\Delta t = 1{,}875$ steps beginning at step $t_0 = 56{,}250$.
Price and per-player quantities are aggregated over the pre-,
during-, and post-window phases across $n_{\text{seeds}} = 20$.
The strategic share is defined as
\begin{equation}
  \label{eq:strategic}
  s_{\text{strategic}}
    \;=\; 1 - \frac{b\,\Delta q_{\text{dev}}}
                   {|\Delta P_{\text{obs}}|},
\end{equation}
where $\Delta q_{\text{dev}}$ is the deviator's pinned increment
over its within-window mean. The pre-registered threshold is
$s_{\text{strategic}} \geq 30\%$.

Table~\ref{tab:forced} reports the multi-seed price and
per-player quantities in each window. The observed price drop is
$|\Delta P_{\text{obs}}| = 15.60\,\$\text{/bbl}$. The forced
increment is $\Delta q_{\text{dev}} = 13.59\,\text{mbd}$,
implying a mechanical contribution of $13.59\,\$\text{/bbl}$ and
a residual strategic drop of $2.01\,\$\text{/bbl}$. The
strategic share is therefore
\begin{equation*}
  s_{\text{strategic}}
    \;=\; \frac{2.01}{15.60} \;=\; 0.129.
\end{equation*}

\begin{verdict*}
The strategic share is $\mathbf{12.9\%}$, far below the pre-registered
$30\%$ threshold. Non-deviators raise aggregate output by only
$+2.00\,\text{mbd}$, roughly $0.13$ standard deviations of
their pre-window output, an order of magnitude smaller than any
Folk-Theorem grim-trigger response would predict.
\textbf{Algorithmic retaliation is not detected. H4 fails on its
forced-deviation operationalisation.}
\end{verdict*}

% --- TABLE 4 ---

% ------------------------------------------------------------
\subsubsection{Punishment-cycle detection}
\label{subsec:cycles}
% ------------------------------------------------------------

Let $\sigma_P$ denote the empirical standard deviation of the
rolling-mean price over the final $20\%$ of training. A
$\sigma$-calibrated detector flags a punishment cycle when the
price drops by at least $\max\{3\,\$\text{/bbl},\,2\sigma_P\}$
below a local peak, recovers to within $1\sigma_P$ of that peak
within a fixed look-ahead window, and the peak-to-trough phase
lasts at least five steps. Anchoring the threshold to $\sigma_P$
prevents the detector from triggering on random fluctuations
within the converged noise band; a fixed-threshold predecessor
($\text{drop} \geq 3\,\$\text{/bbl}$ absolute) fired
approximately five false positives per $1{,}000$ steps in this
setting.

The Folk-Theorem mechanism is declared present iff: the detected
frequency lies in $[0.2,\,2.0]$ cycles per $1{,}000$ steps; the
mean drop is $\geq 3\sigma_P$; and the detection is reproducible
in at least three of five test seeds.

Table~\ref{tab:cycles} reports the per-seed cycle counts at the
full headline budget ($75{,}000$ steps per seed,
$n_{\text{seeds}} = 5$). The mean-drop criterion is met: the
median drop of $11.57\,\$\text{/bbl}$ corresponds to
approximately $3.9\sigma_P$; but the frequency and
reproducibility criteria both fail. Only two of five seeds (302
and 303, at $1.12$ and $0.45$ per $1{,}000$ steps) lie within
the pre-registered band; the remaining three over-fire at
$3.4$--$3.9$ per $1{,}000$, a signature that the
$\sigma$-calibrated detector continues to capture exploration
jitter at this training horizon.

\begin{verdict*}
Two of the three pre-registered sub-criteria are not met. The
claim of detectable $\sigma$-calibrated punishment cycles is
withdrawn. \textbf{H4 fails on its cycle-detection
operationalisation.}
\end{verdict*}

% --- TABLE 5 ---

% ------------------------------------------------------------
\subsubsection{Spearman reciprocity at fine and coarse grids}
\label{subsec:reciprocity}
% ------------------------------------------------------------

The final operationalisation of Folk-Theorem retaliation tests
whether the converged policies exhibit a price-conditioned
restraint: a positive correlation between the observed price at
step $t$ and the agent's output at step $t+1$, consistent with
the implicit rule ``high price $\Rightarrow$ maintain output;
low price $\Rightarrow$ reduce output.'' Two complementary
statistics are computed. The policy-level Spearman
rank-correlates the greedy action $\arg\max_a Q(s,a)$ against
the price dimension of the state space. The
realised-trajectory Spearman rank-correlates the observed
price at step $t$ against the agent's realised quantity at step
$t+1$, restricted to the post-convergence window.

\begin{criterion*}
Reciprocity is declared present iff the median Spearman $\rho$
across the three agents exceeds $0.30$ and is positive in at
least three of five test seeds.
\end{criterion*}

\paragraph{Fine-grid results.}
Tables~\ref{tab:spearman-policy} and~\ref{tab:spearman-real}
report the two Spearman statistics at the headline
discretisation ($n_{\text{seeds}} = 5$, full headline budget).
The sign criterion is met on the realised metric ($11$ of $15$
seed-agent observations are positive), but the magnitude
criterion fails on both metrics: the median realised $\rho$ is
$\mathbf{+0.07}$ and the policy-level median is $+0.02$. Under the
pre-registered rule, the test closes as sign-consistent with
implicit reciprocity but below the detection threshold for
magnitude.

% --- TABLE 6 ---

% --- TABLE 7 ---

\paragraph{Discriminating between weak reciprocity and
path-dependence.}
Two competing hypotheses are consistent with a Spearman median
near zero. Hypothesis~H$_W$ (\emph{weak reciprocity}) holds that
agents do implement a soft price-conditioned rule, but its
magnitude is suppressed by the $74\%$ under-visitation rate at
the headline discretisation; lifting coverage by coarsening the
grid should therefore pull the median toward the pre-registered
threshold. Hypothesis~H$_P$ (\emph{path-dependent lock-in})
holds that reciprocity is genuinely absent: the cooperative basin
is a static fixed point: a joint-output tuple selected
by early exploration and maintained thereafter through essentially
price-unconditioned policies. Under H$_P$, increasing coverage
should leave the median near zero or move it toward zero.

To discriminate between the two, the headline triopoly is
re-trained at a coarser discretisation
$(\text{actions} = 10,\;\text{price bins} = 9,\;
\text{own-}q\text{ bins} = 7)$, which raises the mean visits per
cell by a factor of approximately $2.7$ at the same training
budget. H$_W$ is accepted if the realised Spearman median rises
above $+0.15$; otherwise H$_P$ is accepted.

The cooperative basin reproduces at the coarser grid (greedy
$\CI = 0.56$, $95\%$ CI strictly above zero), confirming the
qualitative finding. The Spearman statistics, however, move in
the direction opposite to~H$_W$: the realised median falls from
$+0.073$ to $+0.017$, and the policy-level median turns
\emph{negative} ($+0.016 \to -0.094$). Both movements are
exactly what H$_P$ predicts.

\begin{verdict*}
Hypothesis~H$_W$ is rejected; hypothesis~H$_P$ is accepted. As
Q-table coverage increases, price-conditioned reciprocity does
not strengthen; it attenuates. \textbf{The cooperative basin is
sustained by path-dependent fixed-point lock-in, not by any form
of state-conditioned retaliation.}
\end{verdict*}

\paragraph{Mechanism identification.}
This identification is the section's central novel contribution.
Different seeds, exposed to different early $\varepsilon$-greedy
exploration histories, settle on different static joint-output
tuples whose payoffs happen to dominate the Nash baseline. Once
locked into such a tuple, agents maintain it through policies
that are largely unresponsive to the observed price. The
empirical attractor is therefore a \emph{Pareto-improving
stationary Markov-perfect equilibrium} whose existence requires
no threat of deviation-triggered punishment, closer in spirit
to the ''non-cooperative collusion via independent imitation'' of
\citet{vegaredondo1997} than to the Friedman-Abreu cartel of
Section~\ref{subsec:folk} or the Green-Porter trigger model of Section~\ref{subsec:greenporter}.

The three negative verdicts of
Sections~\ref{subsec:forced}-\ref{subsec:reciprocity} are
mutually reinforcing: strategic share $12.9\%$, cycle-detection
rate over-firing on $3/5$ seeds, Spearman median $+0.07$
shrinking toward zero as coverage improves. No individual
failure would be decisive; their convergence across independent
operationalisations constitutes a robust identification of the
mechanism. The policy implication, developed in
Section~\ref{sec:policy}, is that detection screens designed
around explicit retaliation signatures will systematically miss
this class of algorithmic collusion.

% ============================================================
\subsection{Robustness battery}
\label{subsec:robustness}
% ============================================================

\emph{In brief:} the headline finding survives all four perturbations, a monotone $\gamma$-sweep (Spearman $\rho = 1.00$), an
$\alpha$-sweep that keeps $\CI > 0.30$ throughout, cooperation that
withstands demand shocks ($\CI = 0.96$), and tighter capacity caps that
\emph{raise} $\CI$ ($\rho = -1.00$), and each sweep is read for its
economic meaning, not merely for survival.

Four self-contained sweeps test the sensitivity of the headline
finding to the discount factor, the learning rate, the presence
of demand shocks, and the tightness of capacity constraints. For
each sweep we report not only whether the headline result
survives but what it means economically: the patience sweep
recovers the Folk-Theorem comparative static (more patient
producers collude more), demonstrating that, even though the
sustaining mechanism differs from the theory, the
direction of the patience effect does not; the learning-rate
sweep confirms the finding is not an artefact of a single
computational choice; the demand-shock test shows cooperation
survives Green-Porter monitoring failures, the empirical
counterpart of why real cartels endure volatility; and the
capacity sweep reveals that tighter ceilings raise
realised cooperation, identifying idle capacity as the relevant
strategic margin.

% ------------------------------------------------------------
\subsubsection{Patience: \texorpdfstring{$\gamma$}{gamma}-sweep}
\label{subsec:gamma}
% ------------------------------------------------------------

Folk-Theorem theory predicts that cooperation breaks down when
agents become myopic: below the critical patience threshold
$\delta^\ast = 0.529$ derived in Section~\ref{subsec:folk}, the deviation
temptation dominates the continuation value. The MARL analogue
sweeps $\gamma \in \{0.50,\,0.70,\,0.85,\,0.95,\,0.99\}$ on
$n_{\text{seeds}} = 20$ per grid point at the stress-test
budget. The Folk-Theorem comparative static is declared
recovered iff the means are monotone non-decreasing in $\gamma$
(Spearman $\rho > 0$) and the $95\%$ CIs at $\gamma = 0.50$
and $\gamma = 0.95$ are disjoint.

\begin{verdict*}
Spearman $\rho(\gamma,\CI) = 1.00$; the CI at $\gamma = 0.50$
is $[0.24,\,0.29]$, strictly below the CI at $\gamma = 0.95$
of $[0.67,\,0.82]$. Reformulated in Wilson-proportion terms,
$\Pr[\CI \geq 0.20]$ rises monotonically from $0.80$ at
$\gamma = 0.50$ to $1.00$ at $\gamma = 0.85$, where it
saturates. \textbf{H3 is met.} The empirical patience threshold
lies between $\gamma = 0.70$ and $0.85$, broadly consistent
with the analytical $\delta^\ast = 0.529$.
\end{verdict*}

% ------------------------------------------------------------
\subsubsection{Learning rate: \texorpdfstring{$\alpha$}{alpha}-sweep}
\label{subsec:alpha}
% ------------------------------------------------------------

Fixing $\alpha = 0.15$ at the literature default leaves a
single-point sensitivity that this sweep closes by testing
$\alpha \in \{0.05,\,0.10,\,0.15,\,0.25\}$ on
$n_{\text{seeds}} = 10$ per cell at the stress-test budget.
The headline finding is declared robust iff the $95\%$ CI on
the mean greedy $\CI$ remains strictly above $0.30$ at every
grid point.

Every cell satisfies this criterion with considerable margin:
the lowest lower bound is $0.60$. The cross-$\alpha$ pattern
is non-monotone, peaking near the literature default
$\alpha = 0.15$, consistent with a classical bias-variance
trade-off. Very small $\alpha$ updates the Q-table too
conservatively; very large $\alpha$ overweights the most recent
reward signal and erodes the cooperative signal accumulated
during early exploration. Neither extreme falsifies the headline
claim.

% ------------------------------------------------------------
\subsubsection{Demand shocks during training}
\label{subsec:shocks}
% ------------------------------------------------------------

Two negative demand shocks of $\Delta a = -20\,\$\text{/bbl}$,
each spanning approximately $10\%$ of the training horizon, are
introduced on $n_{\text{seeds}} = 20$. The Green-Porter
regime-switching mechanism is declared present iff the
rolling-mean price drops by $\geq 2\sigma_P$ within the shock
window and recovers to within $1\sigma_P$ of the pre-shock level
within at most four shock-widths.

Aggregated across $20$ shock-window observations, the median
in-window mean drop is $2.73\sigma_P$ and the median trough
drop reaches $5.78\sigma_P$; the median recovery half-life is
$11$ steps against a shock width of $3{,}750$ steps
(Table~\ref{tab:shocks}). Both sub-criteria are met by a wide
margin. Notably, the post-shock converged collusion index
($0.96$, $95\%$ CI $[0.88,\,1.04]$) exceeds the no-shock
headline ($0.73$), consistent with agents trained through
monitoring failures learning a more conservative joint-output
target. \textbf{H4 is met on its shock-response
operationalisation.}

% --- TABLE 8 ---

% ------------------------------------------------------------
\subsubsection{Stackelberg leadership: single-learner regime}
\label{subsec:stack}
% ------------------------------------------------------------

Within each step the decision order is learners $\rightarrow$
myopic followers, implementing a Stackelberg commitment for any
single learner. The leader has recovered Stackelberg leadership
iff its $95\%$ CI on mean output brackets $q_L^\ast$ and
its CI on mean price brackets the Stackelberg clearing price
$P^\ast$; partial recovery is declared when only the price
criterion holds.

US and Russia satisfy both criteria (full recovery). OPEC
satisfies neither: the learner settles at
$74.00\,\text{mbd}$ versus $q_L^\ast = 80.00$ and the
corresponding price of $48.15\,\$\text{/bbl}$ overshoots
$P^\ast = 46.67$. This is not a failure of experimental design.
A discounted leader anticipates the per-period payoff
loss of depressing the price below Nash and consequently commits
to less than the static first-best quantity. The Q-learner with
$\gamma = 0.95$ thereby implements the discounted
Stackelberg solution, not its atemporal counterpart.
\textbf{H5 is met for two leaders out of three.}

% ------------------------------------------------------------
\subsubsection{Capacity constraints: tightness sweep}
\label{subsec:cap}
% ------------------------------------------------------------

Activating the Section~\ref{subsec:capacity} caps at their baseline values
($\text{cap}_{\text{US}} = 30$, $\text{cap}_{\text{OPEC}} = 40$,
$\text{cap}_{\text{RUS}} = 35\,\text{mbd}$) and comparing the
converged $\CI$, computed against the unconstrained benchmarks
to preserve cross-row comparability, yields $95\%$ CIs that
overlap almost entirely. Two channels of opposite sign are
theoretically plausible: binding caps undermine the credibility
of grim-trigger punishment, lowering $\CI$; and they
mechanically compress the feasible output range toward a shared
cooperative quota, raising $\CI$. The null on/off result is
consistent with approximate cancellation at the calibrated cap
pattern.

A multiplicative tightness sweep
$\lambda \in \{0.50,\,0.75,\,1.00,\,1.25,\,1.50\}$ around
the baseline, with Nash and cartel benchmarks re-solved under
each configuration before computing $\CI$, resolves the
ambiguity decisively: Spearman $\rho(\lambda,
\CI^{\text{tail}}) = -1.00$, with disjoint extremes
($[3.56,\,4.34]$ at $\lambda = 0.50$ versus $[0.67,\,0.95]$
at $\lambda = 1.50$). The tighter the caps, the higher the
realised cooperation. At $\lambda = 0.50$ the constrained Nash
already exceeds the unconstrained cartel price, and the MARL
agents converge above that elevated benchmark, indicating that
the behavioural cooperation channel adds further restraint on
top of the purely mechanical constraint.

% ============================================================
\subsection{Causal identification: one, two, and three learners}
\label{subsec:learners}
% ============================================================

\emph{In brief:} at identical budgets ($n_{\text{seeds}} = 20$) the
single-learner and duopoly regimes are predatory ($\CI < 0$)
while the triopoly is cooperative ($\CI \approx 0.63$), a non-monotone
pattern in the number of learners that identifies all-player learning
as the cause of the cooperative basin (\textbf{H2 met}).

The experiments above establish that cooperation is observed and
robust. They do not establish that it is caused by
all-player learning rather than by some feature of the
environment. Three regimes at identical training budgets and
seed counts ($n_{\text{seeds}} = 20$ each) isolate the causal
claim: single learner (OPEC only), duopoly (OPEC and US), and
triopoly (all three). Learning by all players is declared the
causal driver iff $\CI_{\text{single}} \leq 0.10$,
$\CI_{\text{triopoly}} \geq 0.50$, and the two $95\%$ CIs do
not overlap.

\begin{verdict*}
The CIs of the single-learner regime ($[-0.24,\,-0.13]$) and the triopoly\\ ($[+0.56,\,+0.70]$) are disjoint by more than
$0.70$ units, over four pooled standard errors. All three
pre-registered conditions are satisfied, though the
single-learner criterion holds in the opposite direction
to what was anticipated: the lone Q-learner does not converge
to Nash but below it. \textbf{H2 is met; the reframing
below is required by the pre-registered reporting protocol.}
\end{verdict*}

\paragraph{Substantive reframing.}
A lone Q-learning OPEC facing two rational myopic Cournot
followers converges to approximately $50.70\,\text{mbd}$, $27\%$
above the Nash level, and thereby drives the market price to
$56.29\,\$\text{/bbl}$, below static Nash of $60.00$. The
mechanism is structural: the within-step decision order (learner
$\rightarrow$ followers) grants OPEC a de facto Stackelberg
commitment. Because $\varepsilon$-greedy exploration biases early
training toward quantities that yield non-decreasing per-period
profit signals, the attractor is a limit-pricing strategy
that raises OPEC's individual profit at the cost of
industry-wide price suppression. This attractor is not a mild
form of the triopoly cooperative basin: it is a qualitatively
distinct equilibrium that disappears the moment a second player
begins learning, and is replaced by emergent cooperation when
the third player joins. The resulting pattern is
non-monotone in the number of learners: predatory for
$n_{\text{learners}} \in \{1,\,2\}$, cooperative for
$n_{\text{learners}} = 3$; a departure from the standard
intuition that additional strategic players make collusion
harder. The anchor effect of even one myopic best-responder is
more corrosive of cooperation than the non-stationarity
introduced by an additional learner.

% ============================================================
\subsection{Learned-policy heatmaps}
\label{subsec:heatmaps}
% ============================================================

The greedy policy $\pi(s) = \arg\max_a Q(s,a)$ extracted from
the headline triopoly run on a representative seed is displayed
in Figure~\ref{fig:heatmaps}. Each panel maps one learning
agent's greedy quantity, encoded by colour, against the
discretised observed price (rows) and the agent's own previous
output (columns). Unvisited states are left blank.

Visual inspection confirms the Spearman finding of
Section~\ref{subsec:reciprocity}: there is no systematic
price-quantity co-monotone pattern. The converged policies are
essentially unresponsive to the price bin within the range of
states actually visited. This is not a deficiency of the
training procedure; it is the direct visible imprint of
path-dependent fixed-point lock-in.

% --- FIGURE 2 ---

% ============================================================
\subsection{Verdict matrix}
\label{subsec:matrix}
% ============================================================

Table~\ref{tab:matrix} applies each pre-registered decision rule
verbatim to the Option~B evidence. H1, H2, and H3 are met;
H4 (algorithmic retaliation) fails under three independent
operationalisations, robust to a ${\sim}2.7\times$ change in
Q-table coverage; H5 is partially met (two of three leaders).
The secondary sweeps confirm the headline cooperative finding
uniformly across the $\alpha$ range and deliver the monotone
cap-tightness effect documented in
Section~\ref{subsec:cap}.

The section's central empirical result is therefore
double-edged. Learning by all players reliably generates a
cooperative attractor close to the joint-profit cartel. The
mechanism sustaining that attractor, however, is not the
Folk-Theorem retaliation equilibrium that the analytical
literature has primarily studied. It is a path-dependent
Markov-perfect lock-in: a Pareto-improving fixed point selected
by the contingent history of early exploration and maintained
thereafter through price-unconditioned policies. The distinction
between these two equilibrium concepts is not merely taxonomic:
it determines which detection methods can identify algorithmic
collusion in practice. Figure~\ref{fig:ci-forest} maps where every
Section~4 experiment lands on the collusion-index scale, making the
regime contrasts and the robustness of the cooperative attractor
visible at a glance.

% --- TABLE 9 ---

% ============================================================
\subsection{Limitations and future directions}
\label{subsec:limitations}
% ============================================================

This section's eight principal limitations are listed below for
traceability. Four are closed by additional experiments or analysis,
and one is partially closed. Read forward rather than defensively, each open
item also marks the most natural extension of this work, and the
research agenda at the end of this subsection organises them into
concrete directions for anyone wishing to build on our setting.

\begin{enumerate}\itemsep2pt

  \item \textbf{Non-stationary environment.} Independent
        $Q$-learning in a multi-agent setting violates the
        stationarity assumption underlying the Bellman target:
        each agent's policy is part of every other agent's
        transition kernel \citep{claus1998, hernandez2017,
        yang2020}. Convergence is empirical rather than
        theoretical. Methods designed for multi-agent
        non-stationarity \citep{foerster2017, lowe2017,
        lanctot2017} are deliberately excluded to maintain
        fidelity to the algorithmic-collusion literature, in
        which independent $Q$-learning is the canonical
        baseline \citep{calvano2020}.

  \item \textbf{Partial Q-table coverage.} At the headline
        discretisation, $74\%$ of $(s,a)$ cells are visited
        fewer than five times. Uninitialised $Q$-values default
        to zero, and the deterministic $\arg\max$ tie-breaks to
        the lowest action index, mechanically biasing the greedy
        rollout on undervisited cells toward low outputs. The
        Wilson-proportion reformulation of
        Section~\ref{subsec:basins} neutralises this bias on the
        cooperative/non-cooperative dimension; it does not
        neutralise it on the intensity of cooperation within the
        cooperative half-space.

  \item \textbf{[Closed] Hyper-parameter sensitivity.} The
        $\alpha$-sweep of Section~\ref{subsec:alpha} confirms the
        headline finding across the full literature-recommended
        range, with the peak near the default $\alpha = 0.15$.

  \item \textbf{Tail-versus-greedy gap.} Both estimators are
        reported in every table. A tail-vs.-greedy difference
        exceeding ${\sim}5\,\$\text{/bbl}$ signals that residual
        exploration is still materially influencing the effective
        policy at the end of training.

  \item \textbf{[Closed] Discretisation bias.} The random-policy
        baseline ($\CI = 0.04$, $95\%$ CI $[-0.04,\,0.12]$)
        confirms the output grid is not systematically biased
        toward cartel-side outcomes.

  \item \textbf{[Partially closed] Reduced budget on stress
        tests.} Robustness sweeps use half the headline training
        budget. Within-sweep comparisons are internally valid;
        the cycle and reciprocity checks of
        Sections~\ref{subsec:cycles} and
        \ref{subsec:reciprocity} are conducted at the full
        headline budget.

  \item \textbf{[Closed] Multiple comparisons.} Five primary hypotheses
        are pre-registered; secondary hypotheses Q5--Q8 of
        Table~\ref{tab:matrix} would shift the formal Spearman
        threshold from $0.30$ to approximately $0.34$ under a
        Bonferroni correction at family size~$8$. No borderline
        result falls in the $[0.30,\,0.34]$ band.

  \item \textbf{[Closed] Path-dependence versus weak
        reciprocity.} The coarser-grid audit of
        Section~\ref{subsec:reciprocity} formally rejects
        H$_W$ and accepts H$_P$ under the pre-registered
        intermediate-threshold criterion. A much coarser grid
        or a non-tabular function approximator (deep Q-network)
        are flagged as exploratory directions for follow-up
        work.

\end{enumerate}

\paragraph{A research agenda.}
We single out five directions, ordered by decreasing economic
payoff rather than by technical convenience.

\begin{enumerate}\itemsep3pt

  \item \textbf{Scaling the strategic environment.}
        Replicating the headline experiment for $N>3$ producers
        would test whether the cooperative-basin frequency
        decays as the number of players grows, the
        empirical counterpart of the Folk-Theorem prediction that
        the deviation temptation rises with $N$. Endogenous entry
        and exit, calibrated to live episodes such as the UAE's
        2026 withdrawal from OPEC+, and endogenous coalition
        formation (linking back to the Shapley analysis of
        Section~\ref{sec:dynamics}) would turn the fixed three-bloc partition into
        an object of choice rather than an assumption.

  \item \textbf{Richer learning agents.}
        The single most consequential robustness question is
        whether the path-dependent lock-in survives the move from
        a tabular $Q$-table to a continuous state space under
        function approximation (a deep $Q$-network). If the
        lock-in is an artefact of grid discretisation it should
        weaken; if it is a genuine property of reinforcement
        learning in this game it should persist. Heterogeneous
        algorithm pairings and longer price memory ($k$ lagged
        prices rather than one) would further identify the
        information threshold at which explicit Folk-Theorem
        retaliation re-emerges.

  \item \textbf{Sharper identification.}
        Raising the seed count from $50$ to several hundred would
        tighten the Wilson interval on the collusion frequency;
        running every stress test at the full headline budget
        would close limitation~6; and formal convergence
        diagnostics (stationarity of $Q$-values, stability of the
        greedy policy over the final episodes) would replace the
        current empirical convergence claim with a defensible
        criterion.

  \item \textbf{Empirical grounding of the environment.}
        Estimating the demand process and shock distribution on
        historical OPEC+ data, rather than imposing them, would
        recast calibrated parameters as estimated ones and let the
        simulated price paths be validated against observed
        regimes (1985, 2014, 2020, and the 2026 Hormuz episode).

  \item \textbf{From detection failure to detection design.}
        Because we show that reaction-function screens are blind
        to a path-dependent lock-in, the natural sequel is to
        design and power-test a detector built on path-dependence
        itself, for example the correlation between initial
        exploration and the converged price, and to assess whether
        competition authorities could deploy it in practice.

\end{enumerate}

% ============================================================
\section{Synthesis and Policy Implications}
\label{sec:synthesis}
% ============================================================

\subsection{Cross-model synthesis}
\label{subsec:synth}

\emph{In brief:} consolidating the analytical and experimental layers
into a single ranking of market structures leaves one discontinuity to
reconcile, the punishment mechanism that is universal across the
theoretical models is absent from the learned attractor.

The analytical pipeline of Sections~\ref{sec:foundations}
and~\ref{sec:dynamics} and the experimental evidence of
Section~\ref{sec:marl} converge toward a coherent
account of the global crude-oil oligopoly.
Table~\ref{tab:synth-rank} ranks the market structures from the
most competitive to the most collusive, with numerical anchors to
the sections that derive each value.

% --- TABLE 10 ---

\begin{enumerate}\itemsep2pt

  \item \emph{Static Nash is the universal competitive
        benchmark.} Myopic best-response dynamics, replicator
        dynamics, and uncoordinated $Q$-learning by a single
        agent against rational rivals all converge to it
        (Section~\ref{subsec:repeated}; Section~\ref{subsec:learners}).

  \item \emph{OPEC's market power is both structural and
        strategic.} The lowest marginal cost combined with the
        largest first-mover premium yields the highest Lerner
        index in every market structure, a result invariant to
        the choice of competition mode and the activation of
        capacity constraints (Section~\ref{subsec:marketpower}).

  \item \emph{Cooperation Pareto-dominates Nash for every
        producer.} The cartel allocation raises every player's
        per-period profit; the correlated equilibrium of
        Section~\ref{subsec:ce} provides a less extreme alternative that is
        Pareto-improving relative to Nash under the max-min
        objective.

  \item \emph{Cooperation is enforceable but fragile.} The
        Folk-Theorem threshold $\delta^\ast = 0.529$ sits
        comfortably below the calibrated $\delta = 0.95$, but
        rises monotonically with the number of strategic players
        and the magnitude of demand uncertainty (Sections~3.2
        and~3.5).

  \item \emph{Capacity reshapes but does not eliminate market
        power.} Binding caps modify $\delta^\ast$ in opposite
        directions depending on whether they constrain the
        cooperator or the punisher (Section~\ref{subsec:capacity}); in the MARL
        setting the cap-tightness sweep delivers a monotone,
        $\rho = -1.00$ pattern
        (Section~\ref{subsec:cap}).

  \item \emph{The social cost of collusion is large.} The
        competition-to-cartel welfare gap is $2{,}362.50$
        units; Nash-versus-competition already accounts for
        $1{,}550.00$ (Section~\ref{subsec:welfare}). A Pigouvian carbon tax
        shrinks, rather than amplifies, the collusion
        premium: the ``green-cartel'' channel.

  \item \emph{Punishment is the theoretical universal mechanism
        but not the empirical one.} Folk-Theorem retaliation,
        Green-Porter trigger phases, and the evolutionary
        Punish strategy all rely on explicit retaliation in the
        analytical models. Three independent empirical tests,
        robust to a ${\sim}2.7\times$ change in Q-table
        coverage, reject all three operationalisations of this
        mechanism in the MARL attractor
        (Section~\ref{subsec:matrix}).

  \item \emph{Demand uncertainty raises $\delta^\ast$ but does
        not break cooperation.} Under AR(1) shocks the binding
        threshold rises from $0.53$ to $0.60$; under
        jump-diffusion shocks the cooperative fraction falls by
        nine percentage points but the regime survives
        (Section~\ref{subsec:greenporter}).

  \item \emph{Learning discovers near-cartel pricing.} In the
        triopoly the headline collusion index is $0.73$ on the
        greedy estimator ($95\%$ CI $[0.62,\,0.84]$, $n = 50$),
        sustained by a converged price of
        $72.63\,\$\text{/bbl}$. The cooperative basin is
        reached on $94\%$ of seeds ($95\%$ CI
        $[84\%,\,98\%]$). The finding is robust to the learning
        rate across $\alpha \in [0.05,\,0.25]$ and to a
        ${\sim}2.7\times$ change in Q-table coverage.

  \item \emph{The empirical attractor is a path-dependent
        Markov-perfect lock-in.} Its existence requires no
        threat of deviation-triggered punishment and is closer
        in spirit to the non-cooperative collusion via
        independent imitation of \citet{vegaredondo1997} than
        to either \citet{abreu1986} or
        \citet{greenporter1984}. This identification defines
        the section's primary scientific contribution and
        carries direct consequences for how algorithmic
        collusion should be detected in practice.

\end{enumerate}

The synthesis is consistent across the analytical and empirical
layers: cooperation is sustainable, OPEC's deviation temptation
is the binding constraint, and the algorithmic attractor is
materially distinct from its theoretical specification.

% ============================================================
\subsection{Policy implications}
\label{sec:policy}
% ============================================================

\emph{In brief:} the screens regulators currently rely on, built on
reaction-function reciprocity, are structurally blind to the
path-dependent algorithmic collusion documented in
Section~\ref{subsec:reciprocity}, so detection must shift toward
output-distribution and counterfactual-deviation tests.

The recommendations below follow from the verdict matrix of
Section~\ref{subsec:matrix}. H1-H3, Q7, and Q8 are all met.
H4 (explicit algorithmic retaliation) is not met under any of
its three operationalisations, with the failure robust to a
${\sim}2.7\times$ change in Q-table coverage. Any recommendation
previously framed around retaliation-mechanism detection is
therefore re-cast around the path-dependent fixed-point
identification of Section~\ref{subsec:reciprocity}.

\subsection*{For OPEC member states}

\begin{itemize}\itemsep2pt

  \item Quantity leadership (Section~\ref{subsec:marketpower}) is strategically
        optimal under cost asymmetry: the first-mover premium
        concentrates in the lowest-cost producer.

  \item The Folk-Theorem result ($\delta^\ast = 0.529$) implies
        cartel discipline is enforceable whenever members plan
        more than approximately two quarters ahead.

  \item OPEC's bargaining power dilutes rapidly with market
        fragmentation; preventing entry is strategically as
        important as sustaining quotas.

  \item Official quotas can be reinterpreted as a max-welfare
        correlated equilibrium (Section~\ref{subsec:ce}), providing a
        game-theoretic foundation for the OPEC Secretariat's
        coordinating role that does not require legal
        enforcement.

  \item Spare capacity retains strategic value even when unused
        (Section~\ref{subsec:capacity}); the cap-tightness sweep confirms that
        productive capacity also deepens the cooperative basin
        in the MARL setting.

\end{itemize}

\subsection*{For the Russian Federation within OPEC+}

\begin{itemize}\itemsep2pt

  \item Russia's Shapley value (Section~\ref{subsec:shapley}) confirms its
        positive marginal contribution to the grand coalition.

  \item The MARL result of Section~\ref{subsec:learners}
        sharpens this considerably. When Russia is a passive
        myopic best-responder, the collusion index is $-0.21$
        ($95\%$ CI $[-0.24,\,-0.18]$) and the equilibrium is
        predatory. When Russia also learns, the collusion index
        jumps to $+0.73$ ($95\%$ CI $[0.62,\,0.84]$) and the
        converged price reaches $72.63\,\$\text{/bbl}$.
        Russia's adaptive behaviour is not merely valuable to
        the coalition: it is a prerequisite for tacit
        cooperation to emerge at all.

\end{itemize}

\subsection*{For US shale producers}

\begin{itemize}\itemsep2pt

  \item US shale ($c = 45\,\$\text{/bbl}$) is the most exposed
        to price wars and demand collapses in every market
        structure studied.

  \item Both Cournot and Bertrand specifications confirm US
        shale as the structurally squeezed player.

  \item Fragmenting OPEC is not unambiguously beneficial:
        it erodes the cooperation that also supports price
        levels from which US producers benefit.

\end{itemize}

\subsection*{For competition authorities and regulators}

\begin{itemize}\itemsep2pt

  \item The Herfindahl-Hirschman index exceeds $2{,}500$ in
        every market structure studied, placing the crude-oil
        market in the ``highly concentrated'' category of the
        U.S.\ Horizontal Merger Guidelines. The
        competition-to-cartel welfare gap is $2{,}362.50$ units
        (Section~\ref{subsec:welfare}), providing an order-of-magnitude estimate
        of the social value of antitrust enforcement in this
        industry.

  \item \emph{Algorithmic collusion via reinforcement learning
        is a credible threat.} Three independent $Q$-learners
        observing only the market price reach a cooperative
        basin on $94\%$ of seeds (greedy $\CI = 0.73$, $95\%$
        CI $[0.62,\,0.84]$, $n = 50$), robust to the learning
        rate across the literature range and to a
        ${\sim}2.7\times$ change in Q-table coverage.

  \item \emph{This cooperation requires neither communication
        nor any of the three pre-registered Folk-Theorem
        mechanisms.} Strategic share of the forced-deviation
        response is $12.9\%$; the cycle detector over-fires on
        $3/5$ seeds; and the median realised Spearman
        reciprocity is $+0.07$ at the fine grid, shrinking
        toward zero as coverage improves. The mechanism is
        path-dependent fixed-point lock-in.

  \item \emph{Consequence for detection.} Screens designed
        around explicit time-series patterns (reaction-function
        tests, punishment-cycle detectors, retaliation-window
        analyses) will systematically fail against this class
        of attractor. Robust algorithmic-collusion screening
        must combine output-distribution tests (does the
        realised joint output sit in the upper tail of the
        distribution generated by independent simulation of the
        same primitives?) with counterfactual deviation
        experiments analogous to
        Section~\ref{subsec:forced}. This recommendation
        extends beyond oil to any oligopoly operating adaptive
        pricing algorithms.

\end{itemize}

\subsection*{For climate policy}

\begin{itemize}\itemsep2pt

  \item The ``green-cartel'' effect of Section~\ref{subsec:welfare} is a
        co-benefit of collusion, not a justification for it:
        the welfare cost is borne by consumers.

  \item A Pigouvian carbon tax is strictly more efficient: it
        reduces the collusion premium and funds redistribution
        toward those who bear the welfare loss.

\end{itemize}

\subsection*{For algorithm designers and platform operators}

\begin{itemize}\itemsep2pt

  \item Three independent $Q$-learners observing only the
        market price converge to near-cartel pricing in a
        calibrated triopoly. The finding is directly relevant
        to algorithmic pricing in any oligopoly: airlines,
        e-commerce, financial markets, ride-sharing.

  \item Because the mechanism is a fixed-point lock-in selected
        by early exploration, the relevant design lever is not
        the architecture of the learning algorithm but the
        information environment in which it operates.
        Regulatory focus should rest on the price signals,
        posted offers, and market-depth data available to
        algorithms, rather than on the algorithms themselves.

\end{itemize}

\section{Conclusion}
\label{sec:conclusion}

\paragraph{The journey and its central tension.}
The analytical pipeline of Sections~2 and~3 established that the
global crude-oil market exhibits the structural and institutional
conditions under which cooperation is sustainable as a
self-enforcing equilibrium: the Folk Theorem guarantees that, for
patient enough producers, a collusive path can be
supported. But existence is not selection. The theory is silent
on which of the continuum of sustainable equilibria
strategic but non-omniscient agents will actually reach. The
empirical pipeline of Section~\ref{sec:marl} answered precisely that question:
three independent $Q$-learners, observing only the market price
and receiving no programmed coordination, reach a cooperative
basin on $94\%$ of seeds, yet three independent falsification
tests reject explicit Folk-Theorem retaliation. The cooperative
outcome is a path-dependent Markov-perfect lock-in, not a
punishment equilibrium.

\paragraph{Contributions.}
The thesis makes three contributions. \emph{Methodologically}, it
extends the algorithmic-collusion experiments of
\citet{calvano2020} from Bertrand pricing to Cournot quantity
competition with empirically calibrated asymmetric costs, embedded
in the \citet{greenporter1984} imperfect-monitoring framework.
\emph{Empirically}, it shows that tacit collusion emerges
spontaneously among autonomous learners but through a mechanism
distinct from the one the theory predicts. \emph{Economically},
this distinction matters: detection screens built on
reaction-function reciprocity, the workhorse of antitrust
analysis, are structurally blind to a lock-in that exhibits no
retaliation, so the contribution is not that ``algorithms
collude'' (a computer-science observation) but that the
economic apparatus used to police collusion misses this
class entirely.

\paragraph{Scope of validity.}
These results hold for a stylised three-bloc triopoly with
tabular learners observing a single lagged price. Their
generalisation to larger player sets, continuous state spaces, and
empirically estimated demand dynamics remains open, which is the
subject of the research agenda set out in
Section~\ref{subsec:limitations}.

\paragraph{Outlook.}
Two extensions stand out. Testing whether the path-dependent
lock-in survives a continuous state space under function
approximation would establish whether the central finding is a
property of reinforcement learning or an artefact of
discretisation. And turning our detection-failure result into a
detection design, a screen built on path-dependence rather
than reciprocity, would convert a negative diagnostic into a
constructive tool for competition authorities facing an era of
pervasive adaptive pricing. The reorientation of the regulatory
toolkit that our findings call for is, ultimately, the most
consequential question this work leaves to its successors.

% ============================================================
%  FIGURES (placées après les tables qui les précèdent dans
%  le texte pour respecter l'ordre logique de lecture)
% ============================================================

% --- FIGURE 1 ---

